\section{What the evidence supports, and its limits}\label{ch:limits}

\paragraph{What is established} In generator-based learning with different client models and a long tail, the runs
show three properties of GeFL-F: its generator's rare class rows are shorter, its generated features carry little class
information, and its last layers lose most of what an exact readout recovers (\cref{ch:diagnosis}). The runs show that
FedECS-LT, ending with that readout, has a higher mean than every published method of the setting on both CIFAR
benchmarks and on all six naturally imbalanced datasets, by margins that are large and significant on both CIFAR
benchmarks and on three of the five natural datasets with three seeds; on CIFAR-10-LT the gain is mostly in the recall
of the rare classes (\cref{sec:s-main,sec:s-natural}). They also show that most of this margin is the readout's:
added to the other methods of the GeFL-F pipeline, it leaves FedECS-LT a lead of \R{main/cifar10/er_margin} and
\R{main/cifar100/er_margin} points on the CIFAR benchmarks and at most \R{nat/er_margin_hi} on the natural datasets,
which three seeds do not resolve (\cref{sec:s-same-readout}). The scope of all of this is one client per architecture, a small frozen extractor shared
by all clients, a label shift between clients, and three seeds.

\paragraph{The limits of the study} The paper's limitations section lists what the study does not cover. Each item
limits what may be claimed.
\begin{itemize}
\item \textbf{A small, frozen, shared extractor.} The architectures differ only in their heads. The results say
nothing about pre-trained or larger extractors, or about clients whose whole networks differ. A stronger extractor
would raise every method, and whether it would change the ranking is not known.
\item \textbf{No method built for this exact problem exists.} The comparison is with long-tail losses applied per
client, with FedDF, and with model-heterogeneous methods designed for label skew without a long tail. A margin over
methods built for another problem is weaker evidence than a margin over a direct competitor would be.
\item \textbf{Label shift only.} The splits differ between clients only in their class proportions
(\cref{ch:sums}). The key fact behind the method, that every client's class-$c$ samples point at the same target,
fails under feature shift (the same class looks different on different clients) or concept shift. Those cases are not
tested, and the exactness argument does not cover them.
\item \textbf{One client per architecture} in every reported result. The 50-client runs are pending
(\cref{sec:s-pending}).
\item \textbf{No differential privacy.} The method can be called privacy-preserving only in the sense of secure
aggregation against an honest-but-curious server (\cref{sec:s-privacy}).
\item \textbf{Cost.} FedECS-LT costs more than GeFL-F and far more than per-client long-tail losses
(\cref{tab:s-cost}).
\item \textbf{Low absolute accuracy.} GeFL's own authors attribute GeFL-F's lower accuracy on CIFAR-10 to its small
extractor and heads. In GeFL's balanced setting the paper's earlier implementation of GeFL-F reaches
\R{repro/cifar10/GeFL-F/final}\% on CIFAR-10, close to GeFL's published value of about \R{repro/cifar10/published}\%,
so the low numbers are those of the setting, not of a faulty re-implementation. They are not comparable with
centralised long-tail results for ResNet-32 and larger networks.
\end{itemize}

\paragraph{The trade-offs} Four costs come with the method, and a reader deciding whether to use it should weigh each.
\begin{itemize}
\item \textbf{Each client's own accuracy.} Balanced accuracy and the readout's equal priors serve the federation's
rare classes at the expense of each client's own frequent ones. On the clients' own class mixes FedECS-LT is below
most other methods (\cref{ch:results}).
\item \textbf{Communication and time.} \R{cost/cifar10/comm_ratio_bsm} times the communication and
\R{cost/cifar10/time_ratio_bsm} times the time of Balanced Softmax on CIFAR-10-LT, for a same-readout margin of
\R{main/cifar10/er_diff/Balanced Softmax} points that is not resolved.
\item \textbf{Every client runs every head.} The consensus needs every client to run every architecture's head on the
generated pool each round, and the readout every body once. This is forward computation only, no training, but it
assumes every client can store and run the largest head. The motivation of model heterogeneity, that clients run
different networks because their hardware differs, is therefore only partly kept: the clients still train different
networks, but each must be able to evaluate all of them. The architectures are not secret.
\item \textbf{Released class means.} The server and every client learn the federation's class means and spreads. For a
class held by one client, that is the client's own statistic; for a small class, the mean says more about each image.
With a shallow extractor a class mean is close to a mean image at half resolution.
\end{itemize}

\paragraph{What follows and what does not} It follows that the exact readout raises every method of this setting to which it was
added, by far more than the other differences between them, for one extra round of communication. It follows that
FedECS-LT as a whole has the highest measured balanced accuracy when clients have different models. It does not follow that exact class statistics in the generator are what makes FedECS-LT work: that is
the question of the pending ablations. It does not follow that the decay diagnosis is a cause: that is the question
of the AdamW run. And it does not follow that the result holds with several clients per architecture.

Which questions will a reviewer ask, and what are the honest answers?
